% titel: Blatt 8, A4: Approximation eines Integrals
% skript: Kap. 5 (Riemann-Integral, Approximation)

% Teil I: Werkzeuge, kein Lösungsweg.

Einen numerischen Näherungswert eines Integrals (Gauss-Dichte) aus Python
verstehen und verbessern.

\begin{werkzeugkasten}[Was eine Riemann-Summe wirklich tut]
Der Code
\begin{verbatim}
import math
def f(x): return (1/math.sqrt(2*math.pi))*math.exp(-x**2/2)
S = sum([f(x/100)*(1/100) for x in range(-196, 197)])
\end{verbatim}
zerlegt die Fläche unter dem Graphen in schmale Rechtecke der Breite
$h=\frac1{100}$ an den Stützstellen $x=k/100$, $k=-196,\ldots,196$. Das
ist eine Riemann-Summe für das Integral
\begin{equation*}
  \int_{-1{,}96}^{1{,}96}\frac{1}{\sqrt{2\pi}}\,\mathrm e^{-x^2/2}\,dx .
\end{equation*}
Die Funktion $f$ ist die Dichte der Standard-Normalverteilung
$\mathcal N(0,1)$; das Integral ist die Fläche, die die Dichte zwischen
$-1{,}96$ und $+1{,}96$ einschliesst. Da $1{,}96$ das $97{,}5\%$-Quantil
ist, ist der Flächenanteil zwischen $-1{,}96$ und $+1{,}96$ gerade $0{,}95$.

\textit{Skript:} Kap. 5 (Riemann-Integral).
\end{werkzeugkasten}

\begin{werkzeugkasten}[Warum wird es durch kleinere Schrittweite besser?]
Ein Rechteck unter der Kurve hat auf einem Stück, auf dem die Funktion
glatt ist, einen Fehler proportional zur Schrittweite $h$ (die Kurve
\emph{neigt} sich im Rechteck); über alle Stücke summiert ist der
Gesamtfehler also von Ordnung $O(h)$. Halbiert man $h$, halbiert sich in
etwa der Fehler. Hinzu kommt: \emph{Bessere Quadraturregeln} machen den
Fehler viel schneller klein:
\begin{itemize}
  \item Mittelpunkt-Regel: Stützstelle in der Stückmitte, Fehler
        $O(h^2)$ (die beidseitigen Über-/Unterschätzungen heben sich).
  \item Trapezregel: Verbinden der Höhen durch Geraden, Fehler $O(h^2)$.
  \item Simpson-Regel: Parabeln durch drei Punkte, Fehler $O(h^4)$.
\end{itemize}
Die Aufgabe fragt nach Verbesserung: kleineres $h$ und/oder bessere Regel.

\textit{Skript:} Kap. 5 (numerische Integration).
\end{werkzeugkasten}

\begin{achtung}[Typische Fallen]
\begin{itemize}
  \item Die Breite und den Bereich lesen: $x$ läuft von $-1{,}96$ bis
        $+1{,}96$; ein Rechteck mit linker Stützstelle ist eine
        \emph{Links}-Rechteck-Summe (nicht Mittelpunkt oder Trapez).
  \item Fehlerordnung nicht vergessen: Rechteck $O(h)$, Trapez/Mittelpunkt
        $O(h^2)$, Simpson $O(h^4)$.
  \item Für ein Integral bis $+\infty$ müsste man den Bereich erweitern;
        hier ist $[-1{,}96,1{,}96]$ fest vorgegeben.
\end{itemize}
\end{achtung}
